\begin{ouverture}[college1]{3}{Transformations chimiques et équations}{Partie 1 · Organisation et transformations de la matière}{Comment rendre compte de ce qui disparaît et apparaît ?}
\objectif{Identifier réactifs et produits par des observations et des tests.}
\objectif{Exploiter une équation fournie et compter les atomes.}
\objectif{Utiliser la conservation de la masse avec un système défini.}
\prerequis{Formules moléculaires et coefficients.}
\prerequis{Différence entre physique et chimique.}
\end{ouverture}
\pcPage{Le cours}{Transformations chimiques et équations}
\phantomsection\label{ess-3}
\essSection{1}{Reconnaître une transformation}
Des réactifs sont consommés et des produits se forment. Une couleur, des bulles, un dépôt ou un changement de température sont des indices à interpréter : des bulles peuvent aussi venir d’une ébullition ou d’un gaz dissous. Un test caractéristique, avec réactif et résultat de référence, apporte une information sur l’espèce. Il faut comparer l’état initial et l’état final.\par
\pcFigure{figures/b03-tests.pdf}{Banque de tests : ils sont effectués sous la responsabilité du professeur.}
\essSection{2}{Lire l’équation fournie}
Dans $\ce{CH4 + 2 O2 -> CO2 + 2 H2O}$, le méthane et le dioxygène sont les réactifs ; l’eau et le dioxyde de carbone les produits. Le modèle traduit une redistribution des atomes : un C, quatre H et quatre O sont présents de chaque côté. Les coefficients donnent des proportions d’entités ; ils ne donnent pas directement des masses en grammes.\par
\clearpage
\essSection{3}{Ce qui se conserve}
Au cours d’une réaction chimique, les atomes de chaque élément et la masse totale se conservent dans un bilan matériel complet. Les espèces, le nombre de molécules et les volumes ne sont généralement pas conservés. Dans un système ouvert, une pesée partielle peut varier parce que de la matière entre ou sort. Toujours nommer ce que l’on pèse avant de calculer.\par
\pcFigure{figures/b03-bilan.pdf}{Même masse totale pour un système fermé ; les molécules ont changé.}
\essSection{4}{Construire une conclusion}
Une conclusion relie une donnée au modèle : « l’eau de chaux se trouble ; le test fourni indique du dioxyde de carbone ». Une absence d’effet ne prouve pas toujours une absence totale d’espèce : quantité, sensibilité du test et protocole interviennent. Au Brevet, une démarche partielle pertinente doit être explicitée, avec relation, conversions et unités.\par
\begin{pcEncadre}{method}{Un exemple expliqué}Dans un système fermé, 4 g de dihydrogène et 32 g de dioxygène réagissent entièrement. L’équation fournie est $\ce{2 H2 + O2 -> 2 H2O}$. La masse d’eau formée est $4+32=36\unite g$. On ne multiplie pas 4 g par le coefficient 2 : les masses données concernent déjà les réactifs utilisés.\end{pcEncadre}
\pcPage{Expérimenter et retenir}{Transformations chimiques et équations}
\phantomsection\label{activite-3}
\begin{pcEncadre}{method}{Investigation en classe}Analyser des observations de dissolution, fusion, dégazage et réaction carbonate-acide, avec témoins. Choisir une banque de tests et justifier les conclusions. Le professeur réalise les manipulations autorisées ; prévoir un dispositif adapté au gaz dégagé et ne jamais fermer un récipient rigide improvisé.\end{pcEncadre}
\pcSousTitre{Ce que je dois savoir expliquer}\begin{questions}
\item Définir réactif et produit.
\item Contrôler chaque élément d’une équation fournie.
\item Expliquer une variation de masse apparente en système ouvert.
\end{questions}
\begin{pcEncadre}{attention}{Une erreur à éviter}Une équation se lit avec ses formules : modifier un indice pour faire le compte revient à changer une espèce chimique.\end{pcEncadre}
\pcRepere{cor-activite-3}{Pistes pour l’investigation}\pcGrille
\clearpage\markright{Je démarre}
\rubrique{Je démarre}
\par\noindent\begin{minipage}{\linewidth}
\begin{exercice}[1]{Fondre ou réagir}
\begin{questions}
\item La glace fond. De nouvelles espèces apparaissent-elles ?
\item Le fer rouille. Pourquoi parle-t-on de transformation chimique ?
\end{questions}
\end{exercice}
\end{minipage}
\par\noindent\begin{minipage}{\linewidth}
\begin{exercice}[1]{Avant et après}
\begin{questions}
\item Dans $\ce{C + O2 -> CO2}$, nommer les réactifs.
\item Nommer le produit.
\end{questions}
\end{exercice}
\end{minipage}
\par\noindent\begin{minipage}{\linewidth}
\begin{exercice}[1]{Un test fourni}
\begin{questions}
\item L’eau de chaux se trouble au contact d’un gaz. Quelle espèce est mise en évidence selon le test usuel ?
\item Des bulles seules suffisent-elles ?
\end{questions}
\end{exercice}
\end{minipage}
\par\noindent\begin{minipage}{\linewidth}
\begin{exercice}[1]{Compter}
\begin{questions}
\item Dans $\ce{2 H2 + O2 -> 2 H2O}$, compter H avant et après.
\item Compter O avant et après.
\end{questions}
\end{exercice}
\end{minipage}
\par\noindent\begin{minipage}{\linewidth}
\begin{exercice}[1]{Masse totale}
\begin{questions}
\item Deux réactifs de 12 g et 32 g sont intégralement transformés sans échange de matière. Quelle masse de produits ?
\item Le nombre de molécules doit-il être conservé ?
\end{questions}
\end{exercice}
\end{minipage}
\clearpage\markright{Je m’entraîne}
\rubrique{Je m’entraîne}
\par\noindent\begin{minipage}{\linewidth}
\begin{exercice}[2]{Un gaz s’échappe}
\begin{questions}
\item Un ensemble ouvert pèse 150 g avant et 147,8 g après une réaction. Quelle masse nette a quitté le système si rien d’autre n’entre ni ne sort ?
\item La matière a-t-elle été détruite ?
\end{questions}
\end{exercice}
\end{minipage}
\par\noindent\begin{minipage}{\linewidth}
\begin{exercice}[2]{Une masse inconnue}
\begin{questions}
\item 6 g de carbone réagissent entièrement avec du dioxygène pour former 22 g de $\ce{CO2}$. Trouver la masse de dioxygène consommée.
\item Pourquoi préciser « entièrement » ?
\end{questions}
\end{exercice}
\end{minipage}
\par\noindent\begin{minipage}{\linewidth}
\begin{exercice}[2]{Le bon inventaire}
\begin{questions}
\item Vérifier C, H et O dans $\ce{C2H6O + 3 O2 -> 2 CO2 + 3 H2O}$.
\item Les coefficients 1, 3, 2, 3 sont-ils des rapports de masses ?
\end{questions}
\end{exercice}
\end{minipage}
\par\noindent\begin{minipage}{\linewidth}
\begin{exercice}[2]{Une équation impossible}
\begin{questions}
\item Pourquoi $\ce{H2 + O2 -> H2O}$ ne convient-elle pas comme équation conservant les atomes ?
\item Comparer à $\ce{2 H2 + O2 -> 2 H2O}$.
\end{questions}
\end{exercice}
\end{minipage}
\par\noindent\begin{minipage}{\linewidth}
\begin{exercice}[2]{Des réactifs en excès}
\begin{questions}
\item Un système fermé contient initialement 10 g et 25 g de réactifs. Après, il reste 5 g d’un réactif. Quelle masse totale de produits si aucun autre constituant n’intervient ?
\item Pourquoi ne pas annoncer 35 g de produits ?
\end{questions}
\end{exercice}
\end{minipage}
\clearpage\markright{Je raisonne}
\rubrique{Je raisonne}
\par\noindent\begin{minipage}{\linewidth}
\begin{exercice}[3]{Comparer des contrôles}
\begin{questions}
\item Un réactif de test change de couleur même sans échantillon. Que montre le témoin ?
\item Peut-on alors attribuer la couleur à l’espèce recherchée ?
\end{questions}
\end{exercice}
\end{minipage}
\par\noindent\begin{minipage}{\linewidth}
\begin{exercice}[3]{Mélange et réaction}
\begin{questions}
\item Dissoudre du sucre sans le transformer est-il une réaction ?
\item Pourquoi « mélanger » n’exclut-il pas qu’une réaction se produise dans une autre situation ?
\end{questions}
\end{exercice}
\end{minipage}
\par\noindent\begin{minipage}{\linewidth}
\begin{exercice}[3]{Un solide gagne de la masse}
\begin{questions}
\item Une laine de fer chauffée dans le dioxygène gagne 0,8 g. Quelle matière supplémentaire doit figurer dans le bilan ?
\item Si le fer initial pesait 2,1 g et réagit entièrement, quelle masse finale obtient-on ?
\end{questions}
\end{exercice}
\end{minipage}
\par\noindent\begin{minipage}{\linewidth}
\begin{exercice}[3]{Un protocole de gaz}
\begin{questions}
\item Proposer comment distinguer un gaz dégagé par ébullition et un produit de réaction, sans se fier aux bulles seules.
\item Citer une limite d’une conclusion négative de test.
\end{questions}
\end{exercice}
\end{minipage}
\par\noindent\begin{minipage}{\linewidth}
\begin{exercice}[3]{Prise d’initiative}
\begin{questions}
\item Une réaction donne 18 g d’eau et 44 g d’un gaz. Les seuls réactifs consommés pèsent au total 62 g. Peut-on identifier ce gaz avec les seules masses ?
\item Quel renseignement complémentaire demander ?
\end{questions}
\end{exercice}
\end{minipage}
\pcPage{Défi du chapitre}{Transformations chimiques et équations}
\phantomsection\label{defi-ex-3}
\pcDefiTitre{3}{Le bilan complet}
\begin{questions}
\item Une enceinte contient 8 g de méthane et 40 g de dioxygène. Après réaction, il reste 8 g de dioxygène et il s’est formé 22 g de $\ce{CO2}$. Trouver la masse d’eau formée.
\item Vérifier le bilan matériel.
\item Expliquer pourquoi on ne doit pas compter l’enceinte comme réactif.
\end{questions}
\pcGrille
